\documentclass[12pt,twocolumn]{article}

% Basic packages
\usepackage[utf8]{inputenc}   % UTF-8 encoding
\usepackage{geometry}         % Page margins
\usepackage{amsmath, amssymb} % Math symbols
\usepackage{graphicx}         % Include images
\usepackage{booktabs}         % Nice tables
\usepackage{enumitem}         % Custom lists

% Page setup: 0.5-inch margins
\geometry{
  letterpaper,
  margin=0.5in
}

\usepackage{draftwatermark}
\SetWatermarkText{Do not use in exam}
\SetWatermarkScale{2}        % Size of the watermark (1 = normal)
\SetWatermarkColor[gray]{0.85} % Light gray, can adjust
\SetWatermarkAngle{45}       % Diagonal angle

% START OF DOCUMENT %
\begin{document}

\begin{center}
  {\Large \textbf{Reference Sheet for Exam 2}}
  {\textbf{Johnny A. Rojas }}\\
\end{center}
Hi this is Johnny! The person who wrote this document, due to the restriction
placed on by the professor we are not allowed to use the same personal equation
sheet, I am more than happy for you to copy my cheat sheet so you can make your
own, you can even copy the formating if you like. Sorry for the inconvenience.
\section*{13. Electric Field}
Electric fields can be thought as a shared universal plane that is influenced by
charges and influences charges, similar to gravity and mass.

\subsection*{13.1 Coulomb's Force Law for Point Particles}
Describes the force behaviour that two charges experience.
\begin{align*}
  |\overrightarrow{F}| = F = \frac{1}{4 \pi \epsilon_0} \frac{|Q_1Q_2|}{r^2}
\end{align*}
\subsection*{13.3 Definition of Electric Field}
This equation says that the force on particle 2 is determined by the charge of
the particle 2 and by the electric field $\overrightarrow{E}_1$ made by all
other charged particles.
\begin{align*}
  \overrightarrow{F}_2 = q_2 \overrightarrow{E}_1
\end{align*}
\subsection*{13.4 Electric Field of a Point Particle}
This equation represents how the electric field is influenced due to a charge,
with the use of calculus, it can be used to find the contributions of multiple
electrical fields at a determinent point.
\begin{align*}
  \overrightarrow{E}_1 = \frac{1}{4 \pi \epsilon_0} \frac{q_1}{|\overrightarrow{r}|^2} \hat{r}
\end{align*}
\subsection*{Continuous Charge Distributions}
If we have objects in higher dimensions, the charge is distrbuted along them
continously.
\begin{itemize}
  \item Linear Density (1D, for lines and curves), $\lambda= \frac{dQ}{dl}$ in $(C/m)$.
  \item Surface Density (2D, for planes and surface), $\sigma = \frac{dQ}{dA}$ in $(C/m^2)$
  \item Volume Density (3D, for solids and volumes), $\rho = \frac{dQ}{dV}$ in $(C/m^3)$.
\end{itemize}
Similarly to the total charge for a point is $q$, for the total charge of higher
dimensional objects is the integral of the charge density times the dimensional
measure.
\begin{align*}
  Q = \int \lambda dl \quad Q = \int \sigma dA \quad Q = \int \rho dV
\end{align*}
\subsection*{Electric Field from Continuous Distributions}
Starting from a point-charge field replace $dQ$ with the appropiate dimensional
density.
\begin{align*}
d\overrightarrow{E} = \frac{1}{4 \pi \epsilon_0} \frac{dQ}{r^2} \hat{r}
\rightarrow \overrightarrow{E} = \frac{1}{4 \pi \epsilon_0} \int \frac{dQ}{r^2} \hat{r}
\end{align*}
\subsubsection*{Infinite Line of Charge (uniform $\lambda$)}
\begin{align*}
  E = \frac {\lambda}{2\pi\epsilon_0r}
\end{align*}
\subsubsection*{Ring of a Charge (uniform $\lambda$)}
\begin{align*}
  E = \frac{1}{4\pi\epsilon_0} \frac{Qr}{(r^2+R^2)^{\frac{3}{2}}} \hat{r}
\end{align*}
Note that $R$ is the radious of the ring, $r$ is the distance along the axis
(perpendicular to the plane of the ring) from the center

\end{document}
